The recruitment audits localized the bottleneck; the fixed-mechanism
experiment tested the smallest hypothesized amendment: a local,
novelty-dependent recruitment gain that transiently raises a neuron's
recruitability when its input is unexplained and its post activity is
sparse. The design review \cite{designreview} derived the form from
existing local quantities only --- the allocation rule's unexplained
fraction \(1-R\) and the neuron's own working input current \(I_W\) ---
and concluded that a single new dimensionless magnitude \(k_g\) was
unavoidable (every existing constant is \(\le 1\times\); the required
bootstrap is \(\ge 6\)--\(10\times\) under the measured drive\(\to\)
posts map, whose middle band [4.2, 19] current units is unmeasured by
any committed presentation). Per the pre-registration mandate the
parameter was not fitted: the user froze the first mechanism test at
the lowest defensible value,

\begin{equation}
  i_{\mathrm{boost},i}(t) = k_g \cdot \max(0,\, 1 - R_i(t)) \cdot I_{W,i}(t),
  \qquad k_g = 8.0,
  \label{eq:rg}
\end{equation}

with everything else byte-identical to the committed configuration,
including the identity gate and the frozen endpoint bars
\cite{designreview}. The experiment was executed in two registered
arms.

\subsection{Registration A: uncapped (k\_g = 8.0)}

Nine gain arms (bac/bca/il \(\times\) three seeds) against the nine
committed baselines, all flag-off byte-identical \cite{rg8}.
Endpoints: 4a zero runaway trips (PASS, failures \([]\) in all arms);
4b block-1 first-exposure burst band (FAIL --- pres 1--5 mean
185.9--201.0 Hz vs.\ the [100, 200] Hz bar, with pres-1 at 391--449 Hz
vs.\ 57--94 baseline); 4c first-pattern posts at presentation 20
preserved (PASS, ratio 1.00--1.04); 4d protected cap (PASS, pinned at
0.6000 as in baseline); and the three primary recruitment endpoints
(all FAIL): first-novel posts/n \(\ge 0.5\) (measured 0.00--0.19),
novel permanence \(\ge 150\) (0--22), protected novel end-mass
\(\ge 0.09\) (0.000--0.058).

The mechanism of failure was measured, not inferred: the full-gate
boost on the \emph{first block's own} first exposures
(\(I_W \approx 12\), \(R=0\) \(\Rightarrow\) boost \(\approx 96\),
drive \(\approx 108\) vs.\ the normal 38) drove pres-1 bursts to
391--449 Hz; the hyperactive block then \emph{accelerated} the churn
that destroys the late pattern's substrate (C-cohort prunes in block 1:
195--215 vs.\ 164--187), so at first C exposure the surviving working
current was 0.000--0.14/neuron vs.\ 1.565 baseline --- and the boost is
input-proportional: \(I_W = 0 \Rightarrow\) boost \(= 0\) by
\eqref{eq:rg} itself. The mechanism never engaged at its target moment.
Per the frozen interpretation rule, a 4b failure classifies as a
\emph{recruitment-gain operating-point failure first}; the weak-pattern
recruitment hypothesis was not falsified by registration A
\cite{rg8}.

\subsection{Registration B: threshold-completion cap}

The separately registered amendment \cite{rg8c} clamps the boost at the
current needed to reach threshold in one step,

\begin{equation}
  i_{\mathrm{boost},i}(t) = \min\!\Big(
    k_g \cdot \max(0,1-R_i(t)) \cdot I_{W,i}(t),\;
    \max(0,\,(v_{\mathrm{th}} - v_i(t))\,\tau_m/dt)
  \Big),
\end{equation}

parameter-free (existing \(v_{\mathrm{th}}=1\), \(\tau_m = 20\),
\(dt=1\)); the clamp was implemented per-\emph{neuron} after a
per-synapse first pass was caught in review, with a discriminating unit
test (two afferents, summed raw boost 41.6 \(>\) cap 20 \(\Rightarrow\)
one clamped boost). Identity gate and suites passed; all nine arms
completed, zero failures.

Results against the same bars \cite{rg8c}: 4a PASS; 4b aggregate
5-presentation means 175.6--189.1 Hz inside the band --- but pres-1
remained 378--440 Hz, and under the strict per-presentation reading 4b
is still outside the band in every arm; the cap moved pres-1 only
391 \(\to\) 378 Hz. The mechanism is now understood: the per-tick cap
tops a below-threshold neuron to threshold on every tick while input is
ongoing, sustaining \(\approx 1\) spike/tick (refractory-limited)
regardless of boost magnitude --- the overshoot is
rate-sustained, not magnitude-sustained. The primary endpoints still
failed everywhere (first-novel posts 0.00--0.19; novel permanence
5--15; protected novel end-mass 0.000--0.092, reaching 0.09 in exactly
one arm with only 5 protected synapses); IL preservation stayed partial
(C posts 1.08--1.21\(\times\), permanence 97/74/105 vs.\ the \(\pm20\%\)
of 98 band); gap activity remained elevated (2.3--27.0 vs.\
0.45--5.94 Hz baseline).

\subsection{Why the experiment was closed}

Per the pre-registered outcome rules, each arm was closed at its
failure: no reparameterization, no bracketing, no post-hoc threshold
change, no E-number. The honest verdict language of the records is
preserved here:

\begin{quote}
Registration A: \textbf{operating-point failure} --- 4b fails; the
weak-pattern recruitment hypothesis is \emph{not falsified} (the
mechanism never fired at first exposure because its own block-1 side
effect destroyed the substrate it amplifies).\\[2pt]
Registration B: the cap repairs the aggregate operating point but not
the first-exposure rate; the primary endpoints fail everywhere. The
measured reason is structural: churn precedes exposure (even capped,
block-1 prunes of the late cohort run 193--221 vs.\ 157--187
baseline), and the boost is input-proportional, so a membrane-side gain
has almost nothing to amplify at first contact (surviving working
current 0.00--0.19/neuron vs.\ 1.565). Under the registration B
outcome rule (4b pass + primary fail \(\Rightarrow\) falsified at this
substrate), the records state the hypothesis as \emph{falsified at this
substrate under the corrected operating point}, while noting that
neither registration achieved a clean block-1 operating point, so the
falsification claim is bounded by the residual operating distortion
\cite{rg8c}.
\end{quote}

Both registrations were closed per protocol; the threshold-completion
cap is recorded as the designated, tested alternative; no further
parameter was selected from the results.